\section{Reference Implementation: Issen}
\label{sec:implementation}

This section instantiates Contribution~4 (the layer dependency rule as a
build-system-enforceable constraint). We give the concrete crate-to-layer
mapping of \emph{Issen}, a Rust-based forensic orchestration platform
that wires the five navigation primitives into a unified investigation
pipeline. Issen is not a monolithic tool; it is a constellation of
independently-released crates whose layer membership is fixed by the
taxonomy and whose dependency graph is a direct realization of
Theorems~\ref{thm:parser-indep}~and~\ref{thm:layer-rule}.

\subsection{Architecture}

Table~\ref{tab:crates} summarizes the principal crates and their layers.
The KNOWLEDGE layer is occupied by a single crate, \texttt{forensicnomicon},
which holds zero-dependency, compile-time format constants (page sizes,
EVTX chunk magic, ESE record markers, struct-field offsets). The CONTAINER
layer holds two crates: \texttt{ewf} for E01 disk images and
\texttt{memf-format} for memory dump containers (WinPMEM, raw,
\texttt{hiberfil.sys}, ELF cores). Log containers (EVTX binary, journal
binary, \texttt{tracev3}, PCAP) are absorbed into their respective LOG
FORMAT crates because their outer container and inner format are not
separable.

\begin{table*}[!t]
\centering
\caption{Issen Crate $\to$ Layer Mapping}
\label{tab:crates}
\begin{tabular}{@{}p{0.20\textwidth}p{0.13\textwidth}p{0.06\textwidth}p{0.49\textwidth}@{}}
\toprule
\textbf{Crate} & \textbf{Layer} & \textbf{Type} & \textbf{Role} \\
\midrule
\texttt{forensicnomicon} & KNOWLEDGE & --- & Format constants, magics, struct offsets, schemas \\
\texttt{ewf} & CONTAINER & \src{P} & E01 segment decode $\to$ raw sector stream \\
\texttt{memf-format} & CONTAINER & \src{M} & Dump format decode $\to$ raw page stream \\
\texttt{ext4fs-forensic} & NAVIGATE / FILESYSTEM & \src{P} & Path $\to$ inode $\to$ blocks (ext2/3/4) \\
\texttt{4n6mount} & NAVIGATE / FILESYSTEM & \src{P} & FUSE bridge: image $\to$ OS-visible path \\
\texttt{memf-hw} & NAVIGATE / PAGING & \src{M} & VA $\to$ PA: PML4, PAE, AArch64 page-table walk \\
\texttt{memf-windows} & NAVIGATE / OS STRUCTURE & \src{M} & EPROCESS, VAD, DPAPI, DKOM detection \\
\texttt{winevt-forensic} & NAVIGATE / LOG + PARSER & \src{L} & EVTX chunk seek + BinXML record decode \\
\texttt{issen-remote-access} & NAVIGATE / QUERY & \src{Q} & VQL/WQL/SQL dispatch to remote endpoint \\
\texttt{velociraptor-parser} & NAVIGATE / QUERY & \src{Q} & Decode Velociraptor collection result rows \\
\texttt{cas-forensic} & NAVIGATE / GRAPH & \src{C} & Generic CAS interface (planned) \\
\texttt{git-forensic} & NAVIGATE / GRAPH & \src{C} & Commit/tree/blob DAG traversal (planned) \\
\texttt{sigstore-forensic} & NAVIGATE / GRAPH & \src{C} & Rekor entry resolution (planned) \\
\texttt{browser-forensic} & PARSER & --- & Chromium / Firefox / Safari history; SQLite \\
\texttt{srum-forensic} & PARSER & --- & SRUM / ESE record decode \\
\texttt{issen-correlation} & ORCHESTRATION & --- & Cross-source correlation; Pivot engine \\
\texttt{issen-cli} & ORCHESTRATION & --- & User-facing CLI; routing; presentation \\
\bottomrule
\end{tabular}
\end{table*}

\subsection{Dependency Discipline}

The crate boundaries enforce Theorem~\ref{thm:layer-rule} mechanically.
\texttt{browser-forensic} depends on \texttt{forensicnomicon} (KNOWLEDGE) and
on a minimal SQLite reader; it does \emph{not} depend on \texttt{ewf},
\texttt{ext4fs-forensic}, \texttt{memf-format}, \texttt{memf-hw},
\texttt{memf-windows}, \texttt{winevt-forensic}, or any \src{Q}/\src{C}
crate. Its public API accepts \texttt{\&Path} or \texttt{\&[u8]} -- the
former for live-filesystem and FUSE-mounted access, the latter for any
caller that has already obtained the bytes (memory carving, query result,
git blob retrieval). The five callers in
Section~\ref{sec:cases} all exercise this interface without any modification
to the parser.

The OS STRUCTURE crate \texttt{memf-windows} occupies a notable position. It
locates artifact bytes within a virtual address space (\eg{}, an SQLite page
within \texttt{chrome.exe}'s heap, or a fragment of \texttt{hiberfil.sys}
that contains an EVTX chunk) and then \emph{calls into} a PARSER crate to
interpret those bytes. This is permitted by the dependency rules: NAVIGATE
$\to$ PARSER is a downward dependency (NAVIGATE may use a parser to make
sense of bytes it has located), but PARSER $\to$ NAVIGATE is forbidden.
The asymmetry is what permits parser reuse.

\subsection{Compile-time Parser Registry}

Issen uses Rust's \texttt{inventory} crate to register parsers at link time.
A parser declares itself with a one-line macro invocation:
\begin{quote}\small\ttfamily
inventory::submit!(Parser \{ \\
\hspace*{1em} name: "browser-history", \\
\hspace*{1em} probe: browser\_forensic::probe, \\
\hspace*{1em} parse: browser\_forensic::parse, \\
\hspace*{1em} \ldots \\
\});
\end{quote}
The orchestrator iterates the registry at startup, builds a probe table
keyed by file signatures and content patterns, and routes byte sequences to
parsers without runtime configuration. New parser crates can be added to the
build; the registry picks them up automatically. The architecture fits
Theorem~\ref{thm:parser-indep} naturally: the registry is a set of $\Bstar
\to R$ functions, and the dispatcher cares only about that signature.

\subsection{The PARSER Crate Interface}

PARSER crates expose a uniform interface that admits both source modes:
\begin{quote}\small\ttfamily
pub fn parse\_path(p: \&Path) -\textgreater{} Result\textless{}Vec\textless{}Record\textgreater\textgreater; \\
pub fn parse\_bytes(b: \&[u8]) -\textgreater{} Result\textless{}Vec\textless{}Record\textgreater\textgreater;
\end{quote}
\noindent
The \texttt{parse\_path} variant is a thin convenience over
\texttt{parse\_bytes}: it memory-maps or reads the file and forwards. From
Theorem~\ref{thm:parser-indep}'s standpoint there is only one parser; the
two entry points are presentational sugar for the two most common calling
contexts.

\subsection{Convergence at the Timeline}

Records produced by parsers are normalized into the
\texttt{TimelineEvent} type (and, for derived inferences, into the
\texttt{Evidence} type), and persisted in a DuckDB columnar
store~\cite{raasveldt2019duckdb}. DuckDB serves as the convergence point
for cross-source analytics: a single SQL query can join browser-history
events from a disk-image source (\src{P}), suspicious-process events from
a memory dump (\src{M}), Windows Security log events (\src{L}), live-host
sockstat snapshots (\src{Q}), and signing-log entries (\src{C}). The query
plan does not need to know which navigation primitive produced each row;
that information is preserved as a \texttt{source\_type} column for trust
weighting but does not constrain joins.

\subsection{The \texttt{issen-correlation} Pivot Engine}

The Pivot engine implements the cross-source correlation that the taxonomy
makes principled. A pivot is parameterized by an \emph{anchor} (a value of
some entity type: hash, IP, URL, PID, username) and a \emph{rule set}
(temporal-window joins, spatial joins via filesystem path or memory region,
hash equality joins). Each rule names the source types it draws from. For
example, a ``binary--first-seen'' rule joins \src{P}-source file-creation
events to \src{C}-source Sigstore entries by hash, with a temporal
constraint that the Sigstore entry predates the file creation.

The engine's design exploits Corollary~\ref{cor:c-anchor-free}: \src{C}
matches are weighted with a higher base trust than \src{P}/\src{M}/\src{L}
matches, because the equality of hashes is intrinsic. \src{Q} matches are
weighted by the channel-and-storage trust of the collection.

\subsection{Phases of Implementation}

Issen as of writing realizes \src{P}, \src{M}, and \src{L} fully (the
crates marked without parenthetical ``planned'' in
Table~\ref{tab:crates}). \src{Q} is implemented for VQL via
\texttt{velociraptor-parser} and a remote-access dispatcher. \src{C} is
in active development: \texttt{git-forensic} can traverse a git DAG and
extract commit/tree/blob records; \texttt{sigstore-forensic} ingests Rekor
log exports; the \texttt{cas-forensic} unifying interface is in design. We
report results from \src{P}+\src{M}+\src{L}+\src{Q} in
Section~\ref{sec:cases} as deployed; the \src{C} case study is partially
prospective and we mark its prospective elements explicitly.
